\begin{question}
\noindent A history student is researching the design of an ancient circular observatory, once used to track the positions of stars and planets across the sky. The diagram below shows a circle with centre $O$, representing the circular boundary of the observatory. $AC$ is a diameter of the circle, and $B$ is a point on the circumference.

\begin{center}
\begin{tikzpicture}[scale=2.2]
    % centre and diameter points
    \coordinate (O) at (0,0);
    \coordinate (A) at (180:1);
    \coordinate (C) at (0:1);
    \coordinate (B) at (55:1);

    \draw[name path=circ] (0,0) circle (1);
    \fill[black] (O) circle (0.5pt);

    \draw (A) -- (C);
    \draw (A) -- (B);
    \draw (B) -- (C);

    \node at (A) [anchor=east] {$A$};
    \node at (C) [anchor=west] {$C$};
    \node at (B) [anchor=south] {$B$};
    \node at (O) [anchor=north] {$O$};

    \pic [draw, "$35^\circ$", angle eccentricity=1.5, angle radius=0.5cm] {angle = C--A--B};
    \pic [draw, "$x$", angle eccentricity=1.4, angle radius=0.4cm] {angle = A--C--B};
\end{tikzpicture}
\end{center}

\noindent Angle $BAC = 35^\circ$.

\noindent Work out the size of angle $x$ (angle $ACB$), giving reasons for each stage of your working.
\vspace{4cm}
\begin{flushright}
\answereq{x}{3}
\end{flushright}
\end{question}